\section{Problem setting}\label{sec:setting}

This section fixes the problem class, the SDP that MIST solves, and the linear
system that dominates each interior-point iteration. The rest of the paper is
about solving that linear system quickly and exactly.

\subsection{Multistage polynomial optimization and its SOS relaxation}

We consider polynomial problems whose variables are organized in stages
$\xi_0,\xi_1,\ldots,\xi_N$ and whose objective and constraints couple only
nearby stages:
\begin{equation}\label{eq:pop}
\min_{\xi}\ \sum_{t=0}^{N-1}f_t(\xi_t,\xi_{t+1})
\quad\text{s.t.}\quad
g_{t,a}(\xi_t,\xi_{t+1})\ge0,\quad h_{t,b}(\xi_t,\xi_{t+1})=0,\qquad t=0,\ldots,N-1,
\end{equation}
with polynomial $f_t$, $g_{t,a}$ and $h_{t,b}$. Coupling across a window of
$\ell>2$ stages and a bounded number of global variables are allowed; both enter
the analysis through the parameters of \cref{ass:support}. In robot trajectory
optimization, $\xi_t$ collects the state, control and contact variables of time
step $t$, the equalities encode dynamics, and contact enters through polynomial
complementarity, for example $\lambda\ge0$, $d\ge0$ and $\lambda d=0$, or through
binary mode indicators $\lambda_j^2=\lambda_j$.

Write~\eqref{eq:pop} as $\min\{f(\xi):g_a(\xi)\ge0,\ h_b(\xi)=0\}$. A value $\gamma$
is a certified lower bound if
\begin{equation}\label{eq:sos}
f-\gamma=\sum_c\Bigl(\sigma_{c,0}+\sum_{a}\sigma_{c,a}\,g_a\Bigr)+\sum_b\tau_b\,h_b,
\qquad \sigma_{c,a}=v_{c,a}^{\top}Q_{c,a}v_{c,a},\quad Q_{c,a}\succeq0,
\end{equation}
where $c$ ranges over the cliques of a correlative sparsity
pattern~\cite{waki2006sums,lasserre2006convergent}, $v_{c,a}$ are monomial bases
(possibly pruned by term sparsity~\cite{wang2021tssos,wang2022cstssos}), and
$\tau_b$ are polynomial multipliers. Matching the coefficients of both sides
of~\eqref{eq:sos} gives one affine equation per monomial.

\begin{example}[coefficient matching]\label{ex:matching}
For $p(\xi)=p_0+p_1\xi+p_2\xi^2$ and $v=(1,\xi)$, the identity $p=v^\top Qv$ with
$Q\in\Sym^2_+$ gives three equations, one per monomial $1,\xi,\xi^2$:
$Q_{11}=p_0$, $2Q_{12}=p_1$ and $Q_{22}=p_2$. In a multistage relaxation, a
monomial in the variables of stage $t$ occurs only in the equations of Gram
blocks attached to stages near $t$. This locality is what we exploit.
\end{example}

The lower bound from~\eqref{eq:sos} and the cost of a feasible point recovered
from the relaxation bracket the optimal value of~\eqref{eq:pop}. Whether the
relaxation is tight is a property of the model; the results below concern the cost
of solving a fixed relaxation.

\subsection{The SDP and the interior-point Newton system}

Fixing the relaxation order $r$ and the sparsity rules gives the SDP
\begin{equation}\label{eq:sdp}
\min_{X,\,z}\ \sum_{c=1}^{C}\langle C_c,X_c\rangle+c_z^\top z
\quad\text{s.t.}\quad
\sum_{c=1}^{C}\mathcal A_c(X_c)+Ez=b,\qquad X_c\in\Sym_+^{k_c},\quad z\in\R^{n_f}.
\end{equation}
Each of the $m$ equations is a coefficient-matching equation for one monomial;
$X_c$ are the Gram matrices, of order $k_c$, and $z$ collects the free variables
(the bound $\gamma$ and the coefficients of the multipliers $\tau_b$). We write
$\mathcal A_c(X_c)=(\langle A_i^{(c)},X_c\rangle)_{i=1}^m$.

MIST embeds the free variables in a second-order cone: it adds a scalar $t$ with
zero cost that appears in no equation, and requires $(t,z)\in\mathcal Q^{n_f+1}$.
Every $z$ can be completed by a large enough $t$, so the problem is unchanged, but
all variables now lie in symmetric cones. MIST then applies a primal--dual
interior-point method to the homogeneous self-dual
embedding~\cite{ye1994homogeneous} with Nesterov--Todd
scaling~\cite{nesterov1997self,todd1998nesterov} and Mehrotra
predictor--corrector steps~\cite{mehrotra1992implementation}. Eliminating the
primal and dual directions leaves the normal equations $\mathscr H\,\Delta y=f$ in
the equation multipliers, with
\begin{equation}\label{eq:normal}
\mathscr H=\underbrace{\sum_{c=1}^{C}\mathcal A_c\,\mathcal W_c\,\mathcal A_c^{*}}_{H_{\mathrm{psd}}}
+\widehat E\,W_q^2\,\widehat E^{\top},
\qquad \mathcal W_c(X)=W_cXW_c,\quad \widehat E=[\,0\ \ E\,],
\end{equation}
where $W_c\succ0$ is the scaling of Gram block $c$ and $W_q$ that of the
second-order cone. Entrywise, the first term is
\begin{equation}\label{eq:schur}
(H_{\mathrm{psd}})_{ij}=\sum_{c=1}^{C}\tr\bigl(A_i^{(c)}W_cA_j^{(c)}W_c\bigr).
\end{equation}
The Nesterov--Todd scaling of a second-order cone has the form
$W_q^2=\eta^2(2\bar w\bar w^\top-J)$ with $\eta>0$, $J=\operatorname{diag}(1,-I)$
and $\bar w^\top J\bar w=1$. Because $t$ appears in no equation, the cone's
contribution splits into a sparse part and a single rank-one term:
\begin{equation}\label{eq:operator}
\mathscr H=H_0+uu^\top,\qquad H_0=H_{\mathrm{psd}}+\eta^2EE^\top,\qquad
u=\sqrt2\,\eta\,E\bar w_z,
\end{equation}
where $\bar w_z$ is the trailing part of $\bar w$. The matrix $H_0$ couples two
equations only if they share a Gram block or a free variable. The rank-one term is
dense, which is why it is kept out of every sparsity argument and handled by the
Sherman--Morrison formula (\cref{sec:rankone}). MIST factors a regularized base
\begin{equation}\label{eq:shift}
H_\delta=H_0+\delta I\succ0,\qquad \Hfull=H_\delta+uu^\top,
\end{equation}
with a small shift $\delta\ge0$ chosen by the outer method; every result below
holds verbatim for a diagonal shift.

\paragraph{The linear-algebra task.}
Each IPM iteration assembles $H_\delta$ and $u$, factors $H_\delta$ once, and
solves with a small number $q_{\mathrm{rhs}}$ of right-hand sides: the
homogeneous-embedding column, the predictor, the corrector, and the
Sherman--Morrison column $H_\delta^{-1}u$, so $q_{\mathrm{rhs}}=4$. The corrector
depends on the predictor, so the solves are sequential. Assembling and factoring
$H_\delta$ dominates the cost of an iteration;
\cref{sec:structure,sec:hierarchy,sec:complexity} show how to factor it with work
linear in the horizon, and \cref{sec:algorithm} describes the GPU implementation.
All exactness statements concern this linear system for a fixed shift $\delta$.

\subsection{Notation and computational model}

\Cref{tab:notation} collects the main symbols. We count scalar arithmetic
operations and measure storage in scalar words. Numerical bounds assume that the
symbolic plan and its index maps are available; symbolic analysis is accounted for
separately. A bound called linear in $N$ requires the listed local parameters to
stay bounded as $N$ grows.

\begin{table}[htbp]
\centering
\caption{Main notation.}\label{tab:notation}
\small
\begin{tabular}{@{}lp{0.74\linewidth}@{}}
\toprule
Symbol & Meaning\\
\midrule
$N,\ r$ & horizon and relaxation order\\
$m,\ n_f$ & number of equations and of free variables\\
$\mathcal G=(V,\mathcal E)$ & resource graph on the $m$ equations (\cref{def:resource})\\
$\ell,\ q$ & maximum temporal span of an equation; maximum number of equations first appearing at a stage\\
$n_{\mathrm{var}},\ s$ & variables used by one stage; variables shared across a temporal cut\\
$P$ & monomial interface bound $\binom{s+2r}{2r}$\\
$B,\ g$ & global border equations and $g=|B|$\\
$R_1,\ldots,R_K$ & ordered, disjoint retained interface groups; $p=\max_k|R_k|$\\
$I_j,\ n_j$ & interior connected component and its size\\
$\partial I_j,\ b_j$ & retained neighbors of $I_j$ and their number\\
$d_v,\ w_j$ & later neighbors of interior pivot $v$; maximum later interior neighbors in $I_j$\\
$\Phi_1,\ \Phi_2,\ \mathcal B_2$ & accumulated local-front measures and $\sum_jb_j^2$ (\cref{sec:local})\\
$q_{\mathrm{rhs}}$ & right-hand sides per factorization ($q_{\mathrm{rhs}}=4$ in MIST)\\
\bottomrule
\end{tabular}
\end{table}
